% Source: https://artofproblemsolving.com/community/c3h1348084_problem_1
% Stable identity: aops-topic-1348084-q01
If $n=p_1^{e_1}\cdots p_r^{e_r}$, the number of positive divisors is $(e_1+1)\cdots(e_r+1)$. This equals $3$, a prime, only when $r=1$ and $e_1=2$: so $n=p^2$ for a prime $p$, with divisors $1,p,p^2$.

We need $p^2<2500$, i.e. $p<50$ (and $p^2=2500$ is impossible since $50$ is not prime). The primes below $50$ are
\[
2,3,5,7,11,13,17,19,23,29,31,37,41,43,47,
\]
so there are $\boxed{15}$ such integers. The answer $41$ is wrong: it counts the $49$ squares $2^2,\dots,50^2$ minus a few higher powers, but a square of a composite number such as $36$ has more than $3$ divisors.
